\documentclass[a4paper, 11pt]{article}
\usepackage[utf8]{inputenc}
\usepackage[T1]{fontenc}
\usepackage{lmodern}
\usepackage[french]{babel}
\usepackage{amssymb} % For checkbox symbols
\usepackage{amsmath}
\usepackage{hyperref}

\hypersetup{
    colorlinks=true,
    linkcolor=blue,
    filecolor=magenta,      
    urlcolor=cyan,
}

\title{Exercice : Introduction au RGPD pour BTS SIO1}
\author{}
\date{}

\begin{document}
\maketitle
\thispagestyle{empty}

\section*{Objectifs}
\begin{itemize}
    \item Comprendre les notions de base du RGPD
    \item Identifier des données personnelles
    \item Appliquer les principes fondamentaux de protection des données
    \item Détecter les situations non conformes
\end{itemize}

\section*{Contexte}
Vous travaillez dans une petite entreprise de services informatiques. L'entreprise collecte différentes données sur ses clients et ses employés. Le responsable vous demande de vérifier si certaines pratiques respectent le RGPD.

\section*{Questions}
\begin{enumerate}
    \item \textbf{Qu'est-ce qu'une donnée personnelle ?}
    
    Donnez une définition simple et 3 exemples.
    
    \item \textbf{Parmi les informations suivantes, lesquelles sont des données personnelles ?} (cochez les bonnes réponses)
    
    \begin{itemize}
        \item[$\square$] Nom et prénom
        \item[$\square$] Numéro de dossier interne non lié à une personne
        \item[$\square$] Adresse IP
        \item[$\square$] Adresse postale
        \item[$\square$] Type d'imprimante utilisée dans l'entreprise
        \item[$\square$] Numéro de téléphone mobile
    \end{itemize}
    
    \item \textbf{Citez 3 principes fondamentaux du RGPD.}
    
    (Exemples : minimisation des données...)
    
    \item \textbf{Une entreprise stocke les données de ses clients sans limite de durée. Est-ce conforme au RGPD ? Expliquez pourquoi.}
    
    \item \textbf{Un site web collecte l'adresse email d'un utilisateur pour une newsletter. Quel est l'élément obligatoire avant de pouvoir l'utiliser ?}
    
    \item \textbf{Que signifie le droit à l'oubli ?}
    
    Expliquez en 1 à 2 phrases.
    
    \item \textbf{Dans l'entreprise, le mot de passe d'un salarié est visible en clair dans un fichier Excel. Pourquoi est-ce un problème RGPD ?}
    
\end{enumerate}

\end{document}
